\section{Methodology}\label{sec:methodology}

\subsection{Inference under the physical measure}

The baseline statistical model assumes that the observed underlying follows a geometric Brownian motion under the physical measure,
\begin{equation}
 dS_t=\mu S_t\,dt+\sigma S_t\,dW_t^{P}.
 \label{eq:p_gbm}
\end{equation}
For equally spaced observations with interval $\Delta t$, log returns satisfy
\begin{equation}
 R_i=\log\!\left(\frac{S_{t_i}}{S_{t_{i-1}}}\right)
 \sim
 \mathcal N\!\left[
 \left(\mu-\frac{1}{2}\sigma^2\right)\Delta t,
 \sigma^2\Delta t
 \right].
 \label{eq:return_likelihood}
\end{equation}
The likelihood is the product of the conditional return densities. The simple diffusion is intentionally used as a transparent benchmark: the paper's primary object is the incremental effect of parameter uncertainty, not a claim that constant-volatility GBM is a complete model of crude-oil dynamics.

The baseline prior specification is
\begin{align}
 \mu&\sim\mathcal N(0,1),\\
 \sigma&\sim\operatorname{InvGamma}(2,0.1),\qquad \sigma>0.
 \label{eq:priors}
\end{align}
Prior sensitivity is evaluated separately. Sampling is performed in $(\mu,\eta)$ with $\eta=\log\sigma$ so that positivity is automatic. The transformed log posterior is
\begin{equation}
 \log \pi(\mu,\eta\mid\mathcal D)
 =\log L(\mu,e^\eta;\mathcal D)
 +\log p(\mu)+\log p(e^\eta)+\eta+C,
 \label{eq:posterior}
\end{equation}
where the final $\eta$ is the Jacobian term.

A symmetric random-walk Metropolis proposal is used,
\begin{equation}
 \theta'=\theta+\varepsilon,
 \qquad
 \varepsilon\sim\mathcal N(0,\Sigma_q),
\end{equation}
with $\theta=(\mu,\eta)$. The acceptance probability is
\begin{equation}
 \alpha(\theta,\theta')=
 \min\left\{1,
 \frac{\pi(\theta'\mid\mathcal D)}{\pi(\theta\mid\mathcal D)}
 \right\},
\end{equation}
following the Metropolis--Hastings construction \cite{metropolis1953,hastings1970}. The empirical implementation uses multiple chains and reports Gelman--Rubin $\hat R$ diagnostics for $\mu$ and $\sigma$ rather than treating a single acceptance rate as sufficient evidence of convergence. Effective-sample-size calculations and broader prior sensitivity remain robustness extensions.

\subsection{Risk-neutral valuation}

Historical inference and derivative valuation are kept conceptually and computationally separate. Under the Black--Scholes--Merton benchmark,
\begin{equation}
 dS_t=(r-q)S_t\,dt+\sigma S_t\,dW_t^{Q}.
 \label{eq:q_gbm}
\end{equation}
The historical drift $\mu$ in Equation~\eqref{eq:p_gbm} does not enter Equation~\eqref{eq:q_gbm}. Posterior uncertainty about $\mu$ is therefore relevant for physical forecasting but is not propagated into the no-arbitrage price in the baseline model.

For a discretely monitored arithmetic-average call with fixing dates $t_1,\ldots,t_M$, conditional risk-neutral value is
\begin{equation}
 C^Q(\sigma)=e^{-r\tau}
 \mathbb E_{\sigma}^{Q}\!\left[(A_T-K)^+\mid\mathcal F_t\right],
 \label{eq:conditional_price}
\end{equation}
where $\tau=T-t$. The posterior for $\sigma$ induces a posterior distribution over theoretical prices through the push-forward
\begin{equation}
 p(C\mid\mathcal D)=
 \int
 \delta_{C^Q(\sigma)}(C)
 p(\sigma\mid\mathcal D)\,d\sigma.
 \label{eq:pushforward}
\end{equation}
This distribution quantifies parameter uncertainty in a model price. It is not a predictive distribution of portfolio losses and is not interpreted as VaR or expected shortfall.

\subsection{WTI futures representation}

For the WTI APO application, the state variables are the first-nearby futures fixings in Equation~\eqref{eq:apo_average}. A parsimonious risk-neutral benchmark for a futures price is
\begin{equation}
 dF_u=\sigma F_u\,dW_u^Q,
 \label{eq:futures_q}
\end{equation}
with zero drift under the pricing measure for the appropriately margined futures contract in the stylized setup. The implementation does not replace the entire APO with one fixed CL contract. For each remaining fixing date it identifies the first-nearby contract applicable to that date and initializes the simulated fixing from the corresponding point of the observed futures term structure.

For a valuation date inside the averaging month, simulation starts from Equation~\eqref{eq:partial_fixing}: realized fixings enter deterministically and only the remaining first-nearby settlements are stochastic. This distinction implies that two contracts with the same strike and time to final settlement can have different effective volatility exposure if their realized-average components differ.

The baseline single-factor specification is deliberately interpretable. It can be extended to maturity-specific volatilities, correlated futures-curve factors, stochastic volatility, or mean-reverting commodity models. Those alternatives are treated as robustness extensions because changing the underlying dynamics and changing the treatment of parameter uncertainty are separate questions.

\subsection{Monte Carlo valuation}

Arithmetic-average options generally lack the simple closed form available for geometric averages. We therefore evaluate Equation~\eqref{eq:conditional_price} by Monte Carlo. In the standard Asian benchmark, antithetic sampling and a geometric-average control variate are used following the logic of \cite{kemnavorst1990}. Let
\begin{align}
 Y&=e^{-r\tau}(A_T-K)^+,\\
 X&=e^{-r\tau}(G_T-K)^+,
\end{align}
where $G_T$ is the geometric average and $\mathbb E[X]$ is available analytically. The control-variate estimator is
\begin{equation}
 Y_{cv}=Y-\beta\{X-\mathbb E[X]\},
 \qquad
 \beta=\frac{\operatorname{Cov}(Y,X)}{\operatorname{Var}(X)}.
 \label{eq:control}
\end{equation}
Common random numbers are used when evaluating nearby volatility values for the same contract. This reduces simulation noise in the mapping $\sigma\mapsto C^Q(\sigma)$ and makes the full-Bayes/plugin comparison primarily a comparison of statistical decision rules rather than independent Monte Carlo errors.

\subsection{Competing valuation rules}

Four estimators are compared:
\begin{align}
 \widehat C_{FB}
 &=\mathbb E[C^Q(\sigma)\mid\mathcal D],\\
 \widehat C_{PM}
 &=C^Q(\mathbb E[\sigma\mid\mathcal D]),\\
 \widehat C_{Mode}
 &=C^Q(\widehat\sigma_{Mode}),\\
 \widehat C_{MLE}
 &=C^Q(\widehat\sigma_{MLE}).
 \label{eq:estimators}
\end{align}
The first estimator integrates the pricing map over the posterior. The second collapses the posterior to its mean before pricing. The third uses the marginal posterior mode of $\sigma$, estimated from the sampled volatility distribution, and the fourth is a non-Bayesian likelihood benchmark. The marginal volatility mode is not described as a joint MAP estimator: a joint posterior mode depends on the full parameterization, whereas the implemented comparator is a one-dimensional summary of the marginal posterior for the pricing-relevant parameter.

A useful diagnostic is
\begin{equation}
 \Delta_{FB,PM}
 =\widehat C_{FB}-\widehat C_{PM}.
 \label{eq:fb_pm_gap}
\end{equation}
A second-order expansion gives
\begin{equation}
 \Delta_{FB,PM}
 \approx
 \frac{1}{2}C^{Q\prime\prime}(\bar\sigma)
 \operatorname{Var}(\sigma\mid\mathcal D),
 \qquad
 \bar\sigma=\mathbb E[\sigma\mid\mathcal D].
 \label{eq:taylor_gap}
\end{equation}
Equation~\eqref{eq:taylor_gap} supplies the main mechanism tested in both the synthetic and empirical sections: posterior integration can matter only to the extent that posterior uncertainty interacts with nonlinear exposure to volatility.
